\documentclass[aspectratio=169]{beamer}
\usepackage{spc}
\usepackage{textpos}
\begin{document}

\begin{frame}
  \title{\vspace{-5ex}\darkblue Reproductive Biology\\[2ex]
    \it\large\darkgray Same as Previous Assessment}
  \author{\vspace{-10ex}\darkgray\bf Arni Magnusson}
  \date{\darkgreen SPC PAW Meeting\\[0.5ex]
    24 March 2026, Nouméa}
  \titlepage
  \begin{textblock}{0}(10,0.5)
    \includegraphics[width=4cm]{fish_yellowfin}
  \end{textblock}
\end{frame}

% ______________________________________________________________________________

\begin{frame}{Reproductive potential}
  Email from Nick Davies yesterday:\\[3ex]
  \begin{quotation}\noindent
    Yes, the length-based selectivity feature as used for SKJ2025 can be applied
    to any stock assessment. Note that it differs fundamentally from that of
    SS3, in that it conserves the separability assumption for time-invariant
    selectivity-at-age for the fishing mortality calculations.
  \end{quotation}
\end{frame}

% ______________________________________________________________________________

\begin{frame}
  \vspace{4ex}
  These slides:\\[2ex]
  \begin{center}
    \blue\url{https://github.com/PacificCommunity/ofp-sam-yft-2026-assessment}
  \end{center}
\end{frame}

\end{document}
